\section{Limitations And Future Work}

This paper is deliberately not an end-to-end verification result. That is a
scope choice, not a hidden deficiency. The purpose of M4 is to make the
evidence boundary explicit for a finite social-choice artifact: which claims
are certified under the declared encoding, which semantic target is
Lean-backed, and which stronger correctness obligations remain outside the
paper.

M4 does not prove:

\begin{itemize}[leftmargin=*]
\item Python encoder correctness;
\item \code{\_right\_clauses} correctness;
\item \code{FiniteSchema} profile correctness;
\item SAT assignment semantics as Lean social welfare functions;
\item full semantic-to-CNF correctness;
\item CNF artifact correctness;
\item full selector-free encoding equivalence to Lean \code{MINLIB};
\item a structural proof of the Mask-Shape Collapse Law;
\item a new general Sen theorem;
\item family-scale CNF/LRAT/atlas/repair lifts.
\end{itemize}

The general Sen theorem belongs to the M2 semantic bridge and is not a new M4
claim. M4 uses that established project context as background while keeping
its own claim narrower: a finite audit under the declared M4 encoding, a
Lean-backed semantic target for the central right condition, and explicit
future correctness obligations.

\subsection*{Future Level C Obligations}

The future artifact-correctness layer is Level C. It should be treated as a
set of future assurance obligations, not as an implicit premise of the present
paper. Completing Level C would require:

\begin{itemize}[leftmargin=*]
\item a representation map between Lean profiles and finite schema profiles;
\item a mapping between Lean strict preference predicates and generator
  variables;
\item SAT assignment semantics;
\item a proof that Python right clauses encode the Lean right-condition
  semantics;
\item artifact-level correctness for generated CNF.
\end{itemize}

These obligations would strengthen the project from a claim-boundary audit
over a declared encoding into an artifact-level semantic correctness result.
They are not claimed here.

\subsection*{Future Research Directions}

Future work separates into three tracks. The first is Level C artifact
correctness: connecting the Python and CNF artifact layers to the Lean-backed
semantic target. The second is structural explanation: proving, rather than
only finitely auditing, why the observed collapse pattern holds in the Sen24
case. The third is broader application: testing whether the same
claim-boundary discipline helps other social-choice artifacts or future
AI-mediated governance artifacts. That broader track is downstream relevance,
not a claim that M4 solves AI governance or proves alignment properties.
